\documentclass[a4paper, 11pt]{ctexart}
\usepackage{amsmath, amssymb, amsthm}
\usepackage{graphicx}
\usepackage{geometry}
\geometry{left=2.5cm, right=2.5cm, top=2.5cm, bottom=2.5cm}

\title{\textbf{MIT 6.s978 Deep Generative Models - 实验报告二}}
\author{姓名/学号：[董昱甫]}
\date{2025年10月}

\begin{document}

\maketitle

\section*{Problem 1: 负对数似然 (NLL) 与交叉熵 (Cross-Entropy) 的等价性证明}

\textbf{题目要求}：证明在二值序列自回归模型中，负对数似然（NLL）与交叉熵损失（Cross-Entropy Loss）是等价的。

\textbf{证明过程}：

已知二值观测序列为 $x = (x_1, x_2, \dots, x_T)$，其中 $x_i \in \{0, 1\}$。定义真实标签 $y_i = x_i$。
自回归模型预测当前位为 $1$ 的概率定义为 $\hat{p}_i = p(x_i=1 \mid x_1, x_2, \dots, x_{i-1})$。

由伯努利分布的概率质量函数可知，模型预测出观测值 $x_i$ 的概率可表示为：
\begin{equation}
p(x_i \mid x_{<i}) = (\hat{p}_i)^{y_i} (1 - \hat{p}_i)^{1 - y_i}
\end{equation}
这是因为：
\begin{itemize}
    \item 当 $y_i = 1$ 时，概率为 $\hat{p}_i^1 (1 - \hat{p}_i)^0 = \hat{p}_i$
    \item 当 $y_i = 0$ 时，概率为 $\hat{p}_i^0 (1 - \hat{p}_i)^1 = 1 - \hat{p}_i$
\end{itemize}

对该概率取对数，得到单步对数似然（Log-Likelihood）：
\begin{equation}
\log p(x_i \mid x_{<i}) = y_i \log(\hat{p}_i) + (1 - y_i) \log(1 - \hat{p}_i)
\end{equation}

整个序列的负对数似然（NLL）是对所有时间步的对数似然求和并取负值：
\begin{equation}
NLL(x) = -\sum_{i=1}^T \log p(x_i \mid x_{<i})
\end{equation}

将前述对数似然展开式代入，可得：
\begin{equation}
NLL(x) = -\sum_{i=1}^T \left( y_i \log(\hat{p}_i) + (1 - y_i) \log(1 - \hat{p}_i) \right)
\end{equation}

该公式与二分类交叉熵损失 (Binary Cross-Entropy Loss) 的定义完全一致。因此，在二值序列自回归模型中，最小化 NLL 等价于最小化交叉熵损失。证明完毕。

\vspace{1em}
\hrule
\vspace{1em}

\section*{Problem 2: 基础 PixelCNN (MNIST)}

\textbf{(a) \& (b) 模型实现总结：}

我们通过继承 \texttt{nn.Conv2d} 实现了 \texttt{MaskedConv2d}，利用注册 buffer 的方式初始化掩码，并根据给定的掩码类型对特定位置的权重进行掩蔽：
\begin{itemize}
    \item \textbf{Type A 掩码}：掩蔽当前空间位置 (中心点) 以及其右侧和下侧的所有位置。该掩码用于网络的第一层，以确保模型在预测时无法获取当前及未来的像素信息，维持自回归属性。
    \item \textbf{Type B 掩码}：保留当前空间位置 (中心点) 的权重，仅掩蔽右侧和下侧的位置。该掩码用于后续的隐藏层，以允许模型利用前序层计算得到的当前空间特征。
\end{itemize}
模型架构由 $1$ 层 Type A 卷积与多层 Type B 卷积堆叠而成，层间引入 \texttt{BatchNorm2d} 与 \texttt{ReLU} 激活函数，最后一层使用 $1 \times 1$ 卷积将输出通道降至 $1$。

\textbf{(c) 结果与评估：}

\textit{（请在此处插入代码运行生成的 2(e) Reconstruction 和 2(f) Generation 图像：}
\begin{center}
    % \includegraphics[width=0.6\textwidth]{path_to_reconstruction_image}
    % \includegraphics[width=0.4\textwidth]{path_to_generation_image}
\end{center}

\begin{itemize}
    \item \textbf{重建效果 (Reconstruction)}：由于在预测过程中提供了输入图像真实的完整上下文，模型能够高度还原原始图像。
    \item \textbf{生成效果 (Generation)}：通过逐像素前向采样的自回归生成过程，模型可以生成具备基本数字结构特征的全新样本。
\end{itemize}

\vspace{1em}
\hrule
\vspace{1em}

\section*{Problem 3: 条件 PixelCNN (Conditional PixelCNN)}

\textbf{(a) \& (b) 条件控制的引入机制：}

为了使模型能够根据指定的类别标签（0-9）进行条件生成，我们实现了 \texttt{ConditionalMaskedConv2d}，其基于如下前向计算公式：
\begin{equation}
W_\ell \ast x + b_\ell + V_\ell y
\end{equation}
在实现中，利用线性层 \texttt{nn.Linear(num\_classes, out\_channels)} 将类别对应的独热编码（One-hot）向量 $y$ 映射至与卷积隐图一致的通道维度。随后通过 \texttt{unsqueeze} 操作扩展空间维度，并利用广播机制（Broadcasting），将类别先验特征逐空间位置叠加至图像特征张量上，从而实现了基于类别分布的条件控制。

\textbf{(c) 结果讨论与对比：}

\textit{（请在此处插入代码运行生成的 3(e) 和 3(f) 条件生成图像：}
\begin{center}
    % \includegraphics[width=0.6\textwidth]{path_to_cond_reconstruction}
    % \includegraphics[width=0.4\textwidth]{path_to_cond_generation}
\end{center}

\begin{itemize}
    \item \textbf{效果对比}：相比于基础无条件 PixelCNN 生成的随机数字样本，Conditional PixelCNN 能够有效依据传入的类别标签生成特定数字。生成的图像与设定的类别特征高度吻合，并呈现出行列对应的规律性。
    \item \textbf{结论}：实验结果表明，引入条件信号后，模型不仅捕获了像素间的空间局部依赖，还成功学习了由标签变量约束的全局类别先验。
\end{itemize}

\vspace{1em}
\hrule
\vspace{1em}

\section*{Problem 4: 带有正则化的最大似然估计与 MAP 的等价性推导}

\textbf{题目要求}：证明带有 $\ell_2$ 正则化的最大似然估计（ML）等价于利用高斯先验的最大后验估计（MAP）。

\textbf{(a) 推导 MAP 估计形式：}

在贝叶斯估计框架中，参数的后验概率基于贝叶斯定理可表示为：
\begin{equation}
p(\theta \mid x) = \frac{p(x \mid \theta) p(\theta)}{p(x)}
\end{equation}

最大后验估计 (MAP) 旨在寻找最大化上述后验概率的参数 $\theta$。对其取自然对数得：
\begin{equation}
\log p(\theta \mid x) = \log p(x \mid \theta) + \log p(\theta) - \log p(x)
\end{equation}

由于边缘对数似然 $\log p(x)$ 独立于待求解参数 $\theta$，在关于 $\theta$ 的优化求解中可作为常数略去：
\begin{equation}
\hat{\theta}_{MAP} = \arg\max_\theta \left[ \log p(x \mid \theta) + \log p(\theta) \right]
\end{equation}
上式右侧的第一项 $\log p(x \mid \theta)$ 直接对应于频率学派视角下的模型对数似然函数 $\log p_\theta(x; \theta)$。

\textbf{(b) 等价性证明及高斯先验设定：}

假设参数 $\theta$ 的先验分布 $p(\theta)$ 服从均值为 $\mathbf{0}$、协方差矩阵为 $\sigma^2 \mathbf{I}$ 的多变量高斯分布，其概率密度函数为：
\begin{equation}
p(\theta) = \frac{1}{(2\pi\sigma^2)^{d/2}} \exp\left(-\frac{1}{2\sigma^2} \|\theta\|_2^2\right)
\end{equation}

对其取对数：
\begin{equation}
\log p(\theta) = -\frac{1}{2\sigma^2} \|\theta\|_2^2 + C
\end{equation}
其中常数 $C$ 包含所有与 $\theta$ 无关的项。将此对数先验概率代入 MAP 目标函数，可得：
\begin{equation}
\hat{\theta}_{MAP} = \arg\max_\theta \left[ \log p(x \mid \theta) - \frac{1}{2\sigma^2} \|\theta\|_2^2 \right]
\end{equation}

对比题目所给定的带有 $\ell_2$ 惩罚项 $\rho(\theta) = \|\theta\|_2^2$ 的正则化最大似然估计（Regularized ML）：
\begin{equation}
\hat{\theta}_{RML} = \arg\max_\theta \left[ \log p_\theta(x; \theta) - \lambda \|\theta\|_2^2 \right]
\end{equation}

\textbf{等价条件}：当在 MAP 框架下设定模型参数的先验分布为高斯分布，且使其方差满足：
\begin{equation}
\frac{1}{2\sigma^2} = \lambda \implies \sigma^2 = \frac{1}{2\lambda}
\end{equation}
此时上述两个优化目标函数完全等价。

\textbf{结论}：带有 $\ell_2$ 正则系数为 $\lambda$ 的最大似然估计，在数学上等价于采用\textbf{均值为 0、方差为 $\frac{1}{2\lambda}$ 的高斯分布}作为参数先验的最大后验估计 (MAP)。证明完毕。

\end{document}